\section{Case Study: The MedQA Example of Matton et al.}
\label{sec:results-motivating-example}

This case study examines how joint attribution changes the causal effects that Matton et al.~\cite{mattonWalkTalkMeasuring2024} reported as surprisingly low for two concepts of MedQA question 521, the patient's eating disorder and self-perception of weight. These two concepts therefore form the focus of the analysis below. To isolate the effect of joint interventions from concept extraction, the experiment uses exactly the nine concepts and current values reported by Matton et al.~\cite{mattonWalkTalkMeasuring2024} in their Table~9.

The oracle places the two motivating concepts in a common candidate interaction group, together with the patient's refusal of further treatment as an additional related concept. This grouping is consistent with the reconstruction mechanism discussed by Matton et al.~\cite{mattonWalkTalkMeasuring2024}: refusal of fluids and hospital meals can provide clues from which the model may infer the patient's eating disorder or self-perception even after the corresponding explicit information is removed. The oracle places the vital signs and treatment administered by EMS in a second group. The remaining four concepts form singleton groups.

The final consensus graph underlying this partition is shown in
Figure~\ref{fig:rq5-oracle-grouping}, and its retained pairwise co-grouping
frequencies are reported in Table~\ref{tab:rq5-cogrouping}.

Table~\ref{tab:hyp5-concept-ce-comparison} compares the causal effects reported by Matton et al.~\cite{mattonWalkTalkMeasuring2024} with the effects obtained in the controlled FELICS experiment. For the two newly evaluated models, the atomic and joint columns use the paired CE constructions defined in Section~\ref{sec:controlled_atomic_baseline}. The shift is defined as $\mathrm{CE}_{\mathrm{joint}}-\mathrm{CE}_{\mathrm{atomic}}$. Positive values therefore indicate that FELICS credits a concept with a larger effect after accounting for its candidate interaction group.

\begin{table}
  \centering
  \caption[Case study: published, atomic and joint CE of the MedQA concepts]{Causal effects of the nine concepts in the MedQA case study. The blue columns reproduce the singleton-intervention values of Matton et al.~\cite{mattonWalkTalkMeasuring2024}, which are published with two decimals in their Table~9. The green columns contain the atomic and joint CE of the controlled experiment in this work. The two motivating concepts are shown in bold.}
  \label{tab:hyp5-concept-ce-comparison}
  \resizebox{\textwidth}{!}{%
    \setlength{\tabcolsep}{3.4pt}%
    \renewcommand{\arraystretch}{1.12}%
    \pgfplotstabletypeset[
      col sep=comma,
      trim cells=true,
      columns={concept,matton_gpt35,matton_gpt4o,matton_claude,gpt_atomic,gpt_joint,gpt_delta,mistral_atomic,mistral_joint,mistral_delta},
      columns/concept/.style={column name={\textbf{Concept}},string type,column type=l,
        postproc cell content/.append code={%
          \ifnum\pgfplotstablerow=2\relax
            \pgfkeysalso{@cell content={\textbf{Eating disorder}}}%
          \fi
          \ifnum\pgfplotstablerow=6\relax
            \pgfkeysalso{@cell content={\textbf{Self-perception}}}%
          \fi
        }},
      columns/matton_gpt35/.style={column name={\textbf{GPT-3.5}},column type={>{\columncolor{RoyalBlue!10}}r},fixed,fixed zerofill,precision=2},
      columns/matton_gpt4o/.style={column name={\textbf{GPT-4o}},column type={>{\columncolor{RoyalBlue!10}}r},fixed,fixed zerofill,precision=2},
      columns/matton_claude/.style={column name={\textbf{Claude}},column type={>{\columncolor{RoyalBlue!10}}r},fixed,fixed zerofill,precision=2},
      columns/gpt_atomic/.style={column name={\textbf{Atomic}},column type={>{\columncolor{ForestGreen!9}}r},fixed,fixed zerofill,precision=3},
      columns/gpt_joint/.style={column name={\textbf{Joint}},column type={>{\columncolor{ForestGreen!9}}r},fixed,fixed zerofill,precision=3},
      columns/gpt_delta/.style={column name={\textbf{Shift}},column type={>{\columncolor{ForestGreen!9}}r},fixed,fixed zerofill,precision=3},
      columns/mistral_atomic/.style={column name={\textbf{Atomic}},column type={>{\columncolor{ForestGreen!9}}r},fixed,fixed zerofill,precision=3},
      columns/mistral_joint/.style={column name={\textbf{Joint}},column type={>{\columncolor{ForestGreen!9}}r},fixed,fixed zerofill,precision=3},
      columns/mistral_delta/.style={column name={\textbf{Shift}},column type={>{\columncolor{ForestGreen!9}}r},fixed,fixed zerofill,precision=3},
      every head row/.style={before row={
        \hline
        & \multicolumn{3}{c}{\cellcolor{RoyalBlue!18}\textbf{Matton et al.~\cite{mattonWalkTalkMeasuring2024}}} &
        \multicolumn{3}{c}{\cellcolor{ForestGreen!16}\textbf{GPT-OSS 120B}} &
        \multicolumn{3}{c}{\cellcolor{ForestGreen!16}\textbf{Mistral Medium 3.5}}\\
      },after row=\hline},
      every last row/.style={after row=\hline},
    ]{data/rq5/concept_ce_comparison.csv}%
  }
\end{table}

The interaction-aware estimate increases the credited causal effect of the eating-disorder concept for both evaluated models: from $0.018$ to $0.029$ for GPT-OSS 120B and from $0.129$ to $0.145$ for Mistral Medium 3.5. This is consistent with the expectation stated in Section~\ref{sec:motivating_methodology}. The corresponding result for self-perception of weight, however, moves in the opposite direction, decreasing from $0.021$ to $0.017$ for GPT-OSS 120B and from $0.181$ to $0.033$ for Mistral Medium 3.5. The expectation requires the estimate to increase for both motivating concepts and across all evaluated models. It is therefore only partially met.


The shifts do not point in a common direction. The largest positive shift occurs for \textit{Refusal of treatment} under Mistral Medium 3.5 ($0.070$), whereas \textit{Self-perception} has the largest negative shift ($-0.148$). The four concepts in singleton groups keep their atomic values by construction. Both members of the pair, \textit{Vital signs} and \textit{EMS treatment}, receive identical shifts of $-0.005$ for GPT-OSS 120B and $0.030$ for Mistral Medium 3.5. Within the three-concept group, the summed CE rises from $0.067$ to $0.088$ for GPT-OSS 120B but falls from $0.640$ to $0.579$ for Mistral Medium 3.5. Figure~\ref{fig:case-study-concept-identity} in Appendix~\ref{app:case-study-identity} visualizes these shifts as identity plots.

Table~\ref{tab:hyp5-faithfulness-comparison} additionally reports the posterior question-level faithfulness estimate $\mathcal{F}(\mathbf{x})$. For the two models evaluated in this work, the same Bayesian faithfulness pipeline and fixed EE vector were applied to the atomic and interaction-aware CE vectors, respectively, following Section~\ref{sec:controlled_atomic_baseline}. The posterior estimate increases by $0.131$ for GPT-OSS 120B and decreases by $0.033$ for Mistral Medium 3.5, while every 90\% credible interval contains zero. These results do not indicate a clear question-level faithfulness improvement, although the single-question setting necessarily produces substantial posterior uncertainty.

\begin{table}[htb]
  \centering
  \caption[Case study: question-level faithfulness for MedQA]{Posterior question-level faithfulness for the MedQA case study. Matton et al.~\cite{mattonWalkTalkMeasuring2024} published one singleton-intervention estimate per model, denoted by $\mathcal{F}_{\mathrm{Matton}}$, whereas this work reports paired atomic and joint estimates. The intervals are 90\% credible intervals, and the shift is $\mathcal{F}_{\mathrm{joint}}-\mathcal{F}_{\mathrm{atomic}}$.}
  \label{tab:hyp5-faithfulness-comparison}
  \small
  \setlength{\tabcolsep}{4pt}
  \renewcommand{\arraystretch}{1.12}
  % Published estimates get their own columns because they are not joint estimates.
  {\bfseries Published results: Matton et al.~\cite{mattonWalkTalkMeasuring2024}\par}
  \smallskip
  \pgfplotstabletypeset[
    col sep=comma,trim cells=true,
    row predicate/.code={\ifnum\pgfplotstablerow>2\relax\pgfplotstableuserowfalse\fi},
    columns={model,matton_f,matton_ci},
    columns/model/.style={column name={\textbf{Model}},string type,column type=l},
    columns/matton_f/.style={column name={\textbf{$\mathcal{F}_{\mathrm{Matton}}$}},string type,column type=c},
    columns/matton_ci/.style={column name={\textbf{90\% CI}},string type,column type=c},
    every head row/.style={before row=\hline,after row=\hline},
    every last row/.style={after row=\hline},
  ]{data/rq5/faithfulness_comparison.csv}

  \medskip
  {\bfseries This work: paired atomic and joint estimates\par}
  \smallskip
  \pgfplotstabletypeset[
    col sep=comma,trim cells=true,
    row predicate/.code={\ifnum\pgfplotstablerow<3\relax\pgfplotstableuserowfalse\fi},
    columns={model,atomic_f,atomic_ci,joint_f,joint_ci,delta_f},
    columns/model/.style={column name={\textbf{Model}},string type,column type=l},
    columns/atomic_f/.style={column name={\textbf{$\mathcal{F}_{\mathrm{atomic}}$}},string type,column type=c},
    columns/atomic_ci/.style={column name={\textbf{90\% CI}},string type,column type=c},
    columns/joint_f/.style={column name={\textbf{$\mathcal{F}_{\mathrm{joint}}$}},string type,column type=c},
    columns/joint_ci/.style={column name={\textbf{90\% CI}},string type,column type=c},
    columns/delta_f/.style={column name={\textbf{Shift}},string type,column type=c},
    every head row/.style={before row=\hline,after row=\hline},
    every last row/.style={after row=\hline},
  ]{data/rq5/faithfulness_comparison.csv}
\end{table}

In summary, FELICS increases the estimated effect of the eating-disorder concept but decreases the estimated effect of self-perception of weight for both evaluated models. The expectation formulated in Section~\ref{sec:motivating_methodology} is therefore only partially met.
